% Part: first-order-logic
% Chapter: tableaux
% Section: quantifier-rules

\documentclass[../../../include/open-logic-section]{subfiles}

\begin{document}

\olfileid{fol}{tab}{qrl}

\olsection{પરિમાણક નિયમો}

\subsection{$\lforall$ માટેના નિયમો}

\begin{defish}
\AxiomC{\sFmla{\True}{\lforall[x][!A(x)]}}
\RightLabel{\TRule{\True}{\forall}}
\UnaryInfC{\sFmla{\True}{!A(t)}}
\DisplayProof
\hfill
\AxiomC{\sFmla{\False}{\lforall[x][!A(x)]}}
\RightLabel{\TRule{\False}{\lforall}}
\UnaryInfC{\sFmla{\False}{!A(a)}}
\DisplayProof
\end{defish}

\TRule{\True}{\lforall}માં $t$ બંધ પદ છે (એવું પદ જેમાં કોઈ ચલ નથી).
\TRule{\False}{\lforall}માં $a$ એવો !!a{constant} છે જે
\TRule{\False}{\lforall} નિયમની ઉપરની શાખામાં ક્યાંય આવવો ન જોઈએ.
\TRule{\False}{\forall} અનુમાનના $a$ને આપણે \emph{આઇગનચલ} કહીએ
છીએ.\footnote{ઉપરના નિયમમાં $a$ !!a{constant} હોવા છતાં આપણે
``આઇગનચલ'' શબ્દ
વાપરીએ છીએ. તેના ઐતિહાસિક કારણો છે.}

\subsection{$\lexists$ માટેના નિયમો}

\begin{defish}
\AxiomC{\sFmla{\True}{\lexists[x][!A(x)]}}
\RightLabel{\TRule{\True}{\lexists}}
\UnaryInfC{\sFmla{\True}{!A(a)}}
\DisplayProof
\hfill
\AxiomC{\sFmla{\False}{\lexists[x][!A(x)]}}
\RightLabel{\TRule{\False}{\lexists}}
\UnaryInfC{\sFmla{\False}{!A(t)}}
\DisplayProof
\end{defish}

ફરીથી, $t$ બંધ પદ છે અને $a$ એવો !!a{constant} છે જે
\TRule{\True}{\lexists} નિયમની ઉપરની શાખામાં આવતો નથી.
\TRule{\True}{\lexists} અનુમાનના $a$ને આપણે \emph{આઇગનચલ} કહીએ છીએ.

\TRule{\False}{\lforall} અથવા \TRule{\True}{\lexists} અનુમાનની ઉપરની
શાખામાં આઇગનચલ ન આવે એવી શરતને \emph{આઇગનચલ-શરત} કહે છે.

\begin{explain}
યાદ કરો કે $!A$ એવા !!a{formula} હોય જેમાં !!{variable}~$x$ મુક્ત હોય,
તો આપણે તેને~$!A(x)$ લખીને દર્શાવીએ છીએ. એ જ સંદર્ભમાં $!A(t)$ એ
~$\Subst{!A}{t}{x}$નું સંક્ષિપ્ત રૂપ છે. તેથી
\TRule{\False}{\lexists} નિયમને આ રીતે પણ લખી શકીએ:
\begin{prooftree}
  \AxiomC{\sFmla{\False}{\lexists[x][!A]}}
  \RightLabel{\TRule{\False}{\lexists}}
  \UnaryInfC{\sFmla{\False}{\Subst{!A}{t}{x}}}
\end{prooftree}
નોંધો કે $t$ કદાચ $!A$માં પહેલેથી આવતું હોય; દાખલા તરીકે, $!A$
કદાચ~$\Atom{\Obj P}{t,x}$ હોય. તેથી
$\sFmla{\False}{\lexists[x][\Atom{\Obj P}{t,x}]}$ પરથી
$\sFmla{\False}{\Atom{\Obj P}{t,t}}$નું અનુમાન કરવું એ
\TRule{\False}{\lexists}નો યોગ્ય ઉપયોગ છે. પરંતુ
\TRule{\False}{\lforall} અને~\TRule{\True}{\lexists}ની
આઇગનચલ-શરતો મુજબ !!{constant}~$a$ એ~$!A$માં ન આવવો જોઈએ. તેથી
$\sFmla{\False}{\lforall[x][\Atom{\Obj P}{a,x}]}$ પરથી
$\TRule{\False}{\lforall}$ વાપરીને
$\sFmla{\False}{\Atom{\Obj P}{a,a}}$નું યોગ્ય અનુમાન કરી શકાતું નથી.
\end{explain}

\begin{explain}
\TRule{\True}{\lforall} અને \TRule{\False}{\lexists}માં પદ~$t$ પર કોઈ
નિયંત્રણ નથી. બીજી બાજુ, \TRule{\True}{\lexists} અને
\TRule{\False}{\lforall} નિયમોમાં આઇગનચલ-શરત મુજબ !!{constant}~$a$ એ
સંબંધિત અનુમાનની ઉપરની શાખાઓમાં ક્યાંય આવવો ન જોઈએ. તંત્ર યથાર્થ રહે
તે માટે આ શરત જરૂરી છે. આ શરત વિના
$\lexists[x][\formula{A}(x)] \lif \lforall[x][\formula{A}(x)]$ માટે
નીચેનો બંધ !!{tableau} હોત:
\begin{center}
\begin{tableau}{}
  [\sFmla{\False}{\lexists[x][\formula{A}(x)] \lif \lforall[x][\formula{A}(x)]}, just=\TAss
    [\sFmla{\True}{\lexists[x][\formula{A}(x)]},
      just={\TRule{\False}{\lif}[1]}
      [\sFmla{\False}{\lforall[x][\formula{A}(x)]},
        just={\TRule{\False}{\lif}[1]}
        [\sFmla{\True}{\formula{A}(a)},
          just={\TRule{\True}{\lexists}[2]}
          [\sFmla{\False}{\formula{A}(a)},
            just={\TRule{\False}{\lforall}[3]}, close]
        ]
      ]
    ]
  ]
\end{tableau}
\end{center}
પરંતુ $\lexists[x][\formula{A}(x)] \lif
\lforall[x][\formula{A}(x)]$ પ્રામાણ્ય નથી.
\end{explain}

\end{document}
